\subsection{Probabilistic Query Evaluation}
\label{sec:pqe}

Annotated relations can be used for probabilistic query evaluation using
the so-called intensional approach~\cite{suciu2011probabilistic}. First,
assume a finite set $X$ of variables, each variable $x\in X$ being
assigned a probability $\Pr(x)$. $\Pr$ can be extended to a
probability distribution over valuations over $X$, assuming independence
of variables:
$\Pr(\nu)=\prod_{x\mid\nu(x)=\top}\nu(x)\times\prod_{x\mid\nu(x)=\bot}(1-\nu(x))$.
This in turns extends to a probability
distribution over Boolean functions: if $f\in\mathcal{B}[X]$ is a Boolean
function over variables in~$X$, then $\Pr(f)=\sum_{\nu|f(\nu)=\top} \Pr(\nu)$.

Consider a $\mathcal{B}[X]$-relation $\hat I$ of arity~$k$. For any
subinstance $\hat J$ of~$\hat I$, the characteristic Boolean function of
$\hat J$ within $\hat I$ is: \(\Phi_{\hat I}(\hat J)\defeq\bigwedge_{(u,\alpha)\in
	\hat J}\alpha\land\bigwedge_{(u,\alpha)\in\hat I, (u,\alpha)\notin
	J}\lnot\alpha\). The probability of $\hat J$ is then $\Pr(\Phi_{\hat I}(\hat J))$.
For a query $q$ over $\hat I$ in $\RA_k$, we define the marginal probability that a
tuple $t$ of arity~$k$ appears in
the output of $q$ as $\Pr(t\in q(\hat I))\defeq\sum_{\hat J\subseteq\hat I, t\in
		\sem{J}{q}}\Pr(\hat J)$.

Our semantics for extended relational algebra over $\mathbb{K}$-relations is
compatible with probabilistic query evaluation:

% lax begin Lax392996.ProbabilisticEvaluation
\begin{theoremrep}
	\label{thm:pqe}
	For any finite set of variables~$X$, probability distribution $\Pr$
	over $X$, $\mathcal{B}[X]$-relation~$\hat I$
	and relational
	algebra query~$q$ without aggregation, for any tuple~$t$ with same
	arity as $q$,
	$\Pr(t\in q(\hat I))=\Pr\left(\bigvee_{(t,\alpha)\in \ansem{\hat I}{q}}\alpha\right)$.
\end{theoremrep}
% lax end
\begin{proof}
% lax begin Lax392996Proofs.Bridge.theorem_12
	For a $\mathcal{B}[X]$-relation $\hat I$ and a valuation~$\nu$
	over~$X$, we define $\nu(\hat I)$ to be the (unannotated) relation~$J$
	obtained from $\hat I$ by removing all tuples with annotation $\alpha$
	such as $\alpha(\nu)=\bot$ and keeping all other tuples. In other
	words, $\nu(\hat I)=J\iff\Phi_{\hat I}(\hat J)(\nu)=\top$.

	For a $\mathcal{B}[X]$-relation $\hat I$, a query~$q$ without
	aggregation and a tuple~$t$ with same arity as~$q$, we define
	$\Prov(t\in q(\hat I))$, the \emph{Boolean provenance of $t$ for $q$ on
		$\hat I$}, to be the Boolean function that maps a valuation $\nu$
	to $\top$ if $t\in \sem{\nu(\hat I)}{q}$ and to $\bot$ otherwise.

	We first note that $\Pr(t\in q(\hat I))=\Pr(\Prov(t\in q(\hat I)))$.
	Indeed:
	\begin{align*}
		\Pr(\Prov(t\in q(\hat I)))
		= & \sum_{\nu\mid t\in \sem{v(\hat I)}{q}}\Pr(\nu)
		=\sum_{\hat J\subseteq \hat I,t\in\sem{J}{q}}\sum_{\nu\mid\nu(\hat I)=J}\Pr(\nu)
		=\sum_{\hat J\subseteq \hat I,t\in\sem{J}{q}}\sum_{\nu\mid\nu(\hat I)=J}\Pr(\nu)
	\end{align*}
	while
	\begin{align*}
		\Pr(t\in q(\hat I))=\sum_{\hat J\subseteq\hat I,
			t\in\sem{J}{q}}\Pr(\hat J)
		=\sum_{\hat J\subseteq \hat I,t\in\sem{J}{q}}\Pr[\Phi_{\hat I}(\hat J)]
		=\sum_{\hat J\subseteq \hat I,t\in\sem{J}{q}}\sum_{\nu|\Phi_{\hat I}(\hat J)(\nu)=\top}\Pr(\nu)
		=\sum_{\hat J\subseteq \hat I,t\in\sem{J}{q}}\sum_{\nu|\nu(\hat
			I)=J}\Pr(\nu).
	\end{align*}

	It then suffices to show that
	\[\Prov(t\in q(\hat I))=\bigvee_{(t,\alpha)\in\ansem{\hat
				I}{q}}\alpha\]
	or, in other terms, that
	\[t\in \sem{\nu(\hat I)}{q} \quad\text{if and only if}\quad \exists
		(t,\alpha)\in\ansem{\hat I}{q}, \alpha(\nu)=\top.\]
	We proceed by induction on the structure of $q$.

	\begin{itemize}
		\item \textbf{relation} for any relation label~$R$ in the domain
			of~$\mathcal{D}$,
			\(
			t\in \sem{\nu(\hat I)}{R}
			\iff
			t\in \nu(\hat I)(R)
			\iff
			\exists(t,\alpha)\in\hat I(R)\: \alpha(\nu)=\top
			\iff
			\exists(t,\alpha)\in\ansem{\hat I}{R}, \alpha(\nu)=\top.
			\)

		\item \textbf{projection} for $k\in\bN$, $q\in\RA_{k}$,
			$t_1,\dots, t_n$ of terms of max-index $\leq k$,
			\begin{align*}
				t\in\sem{\nu(\hat I)}{\Pi_{t_1,\dots,t_n}(q)}
				 & \iff
				\exists u\in \sem{\nu(\hat I)}{q}, t=(t_1(u),\dots,t_n(u))                             \\
				 & \iff
				\exists (u,\alpha)\in\ansem{\hat I}{q}, \alpha(\nu)=\top \land t=(t_1(u),\dots,t_n(u)) \\
				 & \iff
				\exists (t,\alpha)\in\ansem{\hat I}{\Pi_{t_1,\dots,t_n}(q)}, \alpha(\nu)=\top.
			\end{align*}
		\item \textbf{selection} for $k\in\bN$, $q\in\RA_k$, and
			$\phi$ a Boolean combination of (in)equality comparisons involving terms
			of max-index $\leq k$,
			\begin{align*}
				t\in\sem{\nu(\hat I)}{\sigma_\phi(q)}
				 & \iff
				t\in\sem{\nu(\hat I)}{q}, \phi(u)                                   \\
				 & \iff
				\exists(t,\alpha)\in\ansem{\hat I}{q}, \alpha(\nu)=\top\land\phi(u) \\
				 & \iff
				\exists(t,\alpha)\in\ansem{\hat I}{\sigma_\phi(q)}, \alpha(\nu)=\top.
			\end{align*}
		\item \textbf{cross product} for $k_1,k_2\in\bN$, $q_1\in\RA_{k_1}$,
			$q_2\in\RA_{k_2}$,
			\begin{align*}
				t\in\sem{\nu(\hat I)}{q_1\times q_2}
				 & \iff
				t=(t_1,t_2), t_1\in\sem{\nu(\hat I)}{q_1}\land t_2\in\sem{\nu(\hat I)}{q_2} \\
				 & \iff
				t=(t_1,t_2),
				\exists (t_1,\alpha)\in\ansem{\hat I}{q_1}, \alpha(\nu)=\top\land
				\exists(t_2,\beta)\in\ansem{\hat I}{q_2}, \beta(\nu)=\top                   \\
				 & \iff
				t=(t_1,t_2),
				\exists (t_1,\alpha,t_2,\beta)\in\ansem{\hat
				I}{q_1}\times\ansem{\hat I}{q_2}, (\alpha\hat\land\beta)(\nu)=\top          \\
				 & \iff
				\exists (t,\gamma)\in\ansem{\hat I}{q_1\times q_2}, \gamma(\nu)=\top.
			\end{align*}
		\item \textbf{multiset sum} for $k\in\bN$, $q_1,q_2\in\RA_k$,
			\begin{align*}
				t\in \sem{\nu(\hat I)}{q_1\uplus q_2} & \iff
				t\in\sem{\nu(\hat I)}{q_1}\uplus\sem{\nu(\hat I)}{q_2}
				\\  & \iff
				(\exists(t,\alpha)\in\ansem{\hat
						I}{q_1}\alpha(\nu)=\top)\lor(\exists(t,\alpha)\in\ansem{\hat
						I}{q_2}\alpha(\nu)=\top)
				\\  & \iff
				\exists(t,\alpha)\in\ansem{\hat
					I}{q_1}\uplus\ansem{\hat I}{q_2}\alpha(\nu)=\top
				\\  & \iff
				\exists(t,\alpha)\in\ansem{\hat
					I}{q_1\uplus q_2}\alpha(\nu)=\top.
			\end{align*}

		\item \textbf{duplicate elimination} for $k\in\bN$,
			$q\in\RA_k$,
			\begin{align*}
				t\in\sem{\nu(\hat I)}{\epsilon(q)}
				\iff
				t\in\sem{\nu(\hat I)}{q}
				\iff
				\exists (t,\alpha)\in \ansem{\hat I}{q}, \alpha(\nu)=\top
				\iff
				\exists (t,\beta)\in\ansem{\hat I}{\epsilon(q)},
				\beta(\nu)=\top
			\end{align*}
			(the last equivalence comes from the fact that a disjunction is
			true if and only if one of the disjunct is true).

		\item \textbf{multiset difference}
			for $k\in\bN$, $q_1,q_2\in\RA_k$,
			\begin{align*}
				t\in\sem{\nu(\hat I)}{q_1-q_2}
				 & \iff
				t\in\sem{\nu(\hat I)}{q_1}\land t\notin\sem{\nu(\hat I)}{q_2} \\
				 & \iff
				(\exists\alpha (t,\alpha)\in\ansem{\hat I}{q_1},
				\alpha(\nu)=\top)\land\lnot
				(\exists \beta (t,\beta)\in\ansem{\hat
				I}{q_2},\beta(\nu)=\top)                                      \\
				 & \iff
				(\exists\alpha (t,\alpha)\in\ansem{\hat I}{q_1},
				\alpha(\nu)=\top)\land\lnot\left(
				\bigvee_{\beta\mid(t,\beta)\in\ansem{\hat I}{q_2}}\beta(\nu)=\top)
				\right)                                                       \\
				 & \iff
				(\exists\alpha (t,\alpha)\in\ansem{\hat I}{q_1},
				\alpha(\nu)\land\lnot\left(
					\bigvee_{\beta\mid(t,\beta)\in\ansem{\hat I}{q_2}}\beta(\nu)=\top)
				\right)=\top                                                  \\
				 & \iff
				\exists(t,\alpha')\in\ansem{\hat I}{q_1-q_2}, \alpha'(\nu)=\top.
				\qedhere
			\end{align*}
	\end{itemize}
% lax end
\end{proof}


The reason why this theorem is about queries without aggregation is
that the result of aggregate queries are tuples whose values include
annotations; for a similar result, we would need to talk about the
distribution of data values, or some summary thereof (such as the
expected value). This is out of scope of this work.

Combining Theorems~\ref{thm:rewriting} and~\ref{thm:pqe}, we obtain:
\begin{corollary}
	\label{cor:pqe}
	For any finite set of variables~$X$, probability distribution $\Pr$
	over $X$, $\mathcal{B}[X]$-relation~$\hat I$
	and relational
	algebra query~$q$ without aggregation, for any tuple~$t$ with same
	arity as $q$, if $\hat q$ is the query rewritten from q by applying the
	rewritten rules (R\ref{r1})--(R\ref{rn}) then
	$\Pr(t\in q(\hat I))=\Pr\left(\bigvee_{(t,\alpha)\in \sem{\hat I}{\hat q}}\alpha\right)$.
\end{corollary}

\begin{example}\label{ex:probability}
	Returning again to Examples \ref{ex:semantics} and \ref{ex:provenance},
	assume we independently assign to each $t_i$ in our
	$\mathit{Personnel}$
	$\mathcal{B}[\set{t_1,\dots,t_7}]$-instance a probability $\Pr(t_i)$
  \revision{with $\Pr(t_1)=0.5$ and $\Pr(t_2)=0.7$}.
	We can then compute the probability of $\qex$ results as\revision{, for
  instance for Nairobi}:
			\(\Pr(t_1\land t_2) = 0.5\cdot0.7=0.35\).
\end{example}
